\documentclass[11pt,reqno]{amsart}

\input{metadata.tex}

\usepackage[margin=1.15in]{geometry}
\usepackage{amsmath,amssymb,amsthm}
\usepackage{enumitem}
\usepackage{microtype}
\usepackage{xcolor}
\definecolor{linknavy}{RGB}{20,50,110}
\usepackage[colorlinks=true,linkcolor=linknavy,citecolor=linknavy,urlcolor=linknavy]{hyperref}
\hypersetup{
  pdftitle={Necessary Ambiguity: The Galois Floor and the Orientation Dyad of Jacobian Counterexamples},
  pdfauthor={\AuthorName},
  pdfkeywords={Jacobian Conjecture, Keller map, Galois closure, monodromy, sign character, double cover, orientation}
}

\numberwithin{equation}{section}
\newtheorem{theorem}{Theorem}[section]
\newtheorem{lemma}[theorem]{Lemma}
\newtheorem{proposition}[theorem]{Proposition}
\newtheorem{corollary}[theorem]{Corollary}
\theoremstyle{remark}
\newtheorem{remark}[theorem]{Remark}
\newtheorem{question}[theorem]{Question}

\newcommand{\C}{\mathbb C}
\newcommand{\Z}{\mathbb Z}
\newcommand{\R}{\mathbb R}
\newcommand{\eps}{\varepsilon}
\DeclareMathOperator{\Mon}{Mon}
\newcommand{\SF}{S_F}

\title[Necessary ambiguity: the Galois floor and the orientation dyad]
{Necessary Ambiguity:\\ The Galois Floor and the Orientation Dyad\\ of Jacobian Counterexamples}

\author{\AuthorName}
\address{\AuthorAffiliation}
\email{\AuthorEmail}
\date{\PreprintDate}
\thanks{Version \PreprintVersion. Research note; not peer reviewed. Zenodo DOI: \texttt{\PreprintDOI}.}

\subjclass[2020]{Primary 14R15; Secondary 14D05, 12F10}
\keywords{Jacobian Conjecture, Keller map, Galois closure, monodromy,
nonproperness set, sign character, orientation double cover}

\begin{document}

\begin{abstract}
Using the classical Galois case of the Jacobian problem (Campbell 1973;
Razar; Wright), we record what every counterexample to the Jacobian
Conjecture must carry on its inverse side. (A)~Every Keller counterexample,
in every dimension, has generic fiber degree at least three. (B)~At the
minimum degree the Galois closure is the full symmetric group $S_3$: the
minimal inverse ambiguity is necessarily a complete symmetric triad.
(C)~Consequently every minimal counterexample carries a nontrivial sign
character and its orientation double cover --- a necessary $\Z/2$ of
structure we call the dyad. (D)~A short group-theoretic analysis shows the
dyad is the unique normal intermediate stage of any resolution of the
ambiguity, that no orientation of the ambiguity exists over the base, and
that the two cyclic orientation classes become invariantly distinguishable
exactly on the double cover, where they are exchanged by the deck
involution and (for real-coefficient maps) by complex conjugation: in a
slogan, orientation data is born exactly where the dyad dies. (E)~All of
this is verified in exact arithmetic on the July 2026 counterexample and on
$400$ members of its deformation family; the scripts accompany this
deposit. The results are unconditional theorems of the complex affine
category; no claim outside mathematics is made.
\end{abstract}

\maketitle

\section{Introduction}\label{sec:intro}

Keller's Jacobian Conjecture (1939) asserted that every polynomial map
$F\colon\C^n\to\C^n$ with nonzero constant Jacobian determinant (a
\emph{Keller map}) is a polynomial automorphism \cite{Keller,vdE}. It was
refuted on July 20, 2026 by an explicit map $\C^3\to\C^3$
\cite{Alpoge,Counterexample2026}. This note reverses the direction of
interrogation: rather than asking how a counterexample was constructed, we
ask what \emph{every} counterexample must carry.

The answers assemble from classical results, principally the Galois case of
the problem, proved by Campbell \cite{Campbell} over $\C$ and in general by
Razar \cite{Razar} and, independently, Wright \cite{Wright}:

\begin{quote}
\emph{A polynomial map from $\C^n$ to $\C^n$ has a polynomial inverse if
(and only if) its Jacobian determinant never vanishes and it induces a
normal extension of function fields} \cite[p.~244]{Campbell}.
\end{quote}

\noindent In characteristic zero, normal means Galois
(\cite[Lemma~1]{Campbell}). Our main results:

\begin{enumerate}[label=\textup{(\Alph*)}]
\item every Keller counterexample has generic fiber degree $d\ge3$
(Theorem~\ref{thm:floor});
\item if $d=3$, the monodromy of the inverse cover is the full symmetric
group $S_3$ (Theorem~\ref{thm:floor});
\item hence every minimal counterexample carries a nontrivial sign
character and a canonical orientation double cover
(Section~\ref{sec:dyad});
\item the coupling theorems of Section~\ref{sec:coupling}: the dyad is the
unique normal intermediate of any resolution; no orientation exists over
the base; the two cyclic orientation classes separate exactly on the double
cover and are exchanged there by two parity operations;
\item exact-arithmetic verification on the known counterexample and on
$400$ members of its deformation family (Section~\ref{sec:verification}).
\end{enumerate}

Everything is a theorem of the complex affine category. We note the scope
honestly: over $\R$ the landscape is entirely different --- Pinchuk
\cite{Pinchuk} constructed a real polynomial map $\R^2\to\R^2$ with
nowhere-vanishing Jacobian that is not injective, and the restriction of
the 2026 map to $\R^3$ (its coefficients are rational) is a real Keller
counterexample whose target chambers into regions of three and of one real
preimage, with no monodromy connecting them. The braided structure studied
here is a complex phenomenon; a companion note takes up that contrast.

\section{The covering over the complement of the Jelonek
set}\label{sec:setting}

Let $F\colon\C^n\to\C^n$ be a Keller map; then $F$ is \'etale and dominant,
the function-field extension $\C(x)/\C(F)$ is finite separable, and its
degree $d$ equals the number of preimages of a general point
(\cite[Introduction]{Campbell}). Let $\SF$ denote the set of points at
which $F$ is not proper (the Jelonek set \cite{Jelonek}), a closed
algebraic subset, of pure codimension one and uniruled when nonempty
\cite{Jelonek}. Since $F$ is \'etale, every finite fiber point has local
multiplicity one; by the properness criterion (finite fiber whose
multiplicities sum to the generic degree; see \cite{Jelonek}), $F$ is
proper exactly over
\[
U:=\C^n\smallsetminus \SF ,
\]
and the restriction $\pi\colon F^{-1}(U)\to U$ is a covering map of degree
$d$. Its total space $F^{-1}(U)=\C^n\smallsetminus F^{-1}(\SF)$ is the
complement of a proper analytic subset of $\C^n$, hence connected;
therefore the monodromy $\Mon(\pi)\le S_d$ is \emph{transitive}, and it is
isomorphic to the Galois group of the Galois closure of $\C(x)/\C(F)$.

\begin{lemma}\label{lem:horizon}
If $F$ is a Keller counterexample, then $\SF\neq\emptyset$.
\end{lemma}

\begin{proof}
If $\SF=\emptyset$ then $F$ is proper, and properness is one of Campbell's
sufficient conditions for a polynomial inverse
(\cite[Introduction, condition 2]{Campbell}); directly: a proper \'etale
map onto the simply connected $\C^n$ is a finite covering, hence trivial,
hence $d=1$ and $F$ is birational, and Keller's theorem \cite{Keller}
applies.
\end{proof}

\section{The floor}\label{sec:floor}

\begin{theorem}[The Galois floor]\label{thm:floor}
Let $F$ be a Keller counterexample of generic degree $d$. Then:
\begin{enumerate}[label=\textup{(\roman*)}]
\item $d\neq1$ \textup{(}Keller \cite{Keller}: the birational case\textup{)};
\item $d\neq2$ \textup{(}Campbell's Corollary \cite[p.~248]{Campbell}: the
degree of a nowhere degenerate polynomial map cannot be two, since
quadratic extensions are normal\textup{)};
\item hence $d\ge3$; and
\item if $d=3$, then $\Mon(\pi)=S_3$.
\end{enumerate}
\end{theorem}

\begin{proof}
(i) and (ii) are the cited classical results. For (iv): the transitive
subgroups of $S_3$ are $A_3$ and $S_3$. If $\Mon(\pi)=A_3$, its order
equals the degree, so the Galois closure of $\C(x)/\C(F)$ has degree
$3=[\C(x):\C(F)]$; the extension is its own closure, hence normal, and
Campbell's Theorem makes $F$ invertible --- a contradiction.
\end{proof}

\begin{remark}[Provenance]\label{rem:provenance-floor}
The floor stands on three named results --- Keller (degree one), Campbell's
Corollary (degree two), Campbell's Theorem (cyclic degree three) --- and
the assembly is one line. Its interest is entirely occasioned by
\cite{Alpoge}: for eighty-seven years the class it quantifies over was not
known to be nonempty. To our knowledge the consequence
(iv) and the structural analysis below have not appeared in print; we make
no priority claim pending specialist review.
\end{remark}

\begin{proposition}[Descent]\label{prop:descent}
For any counterexample, $\Mon(\pi)$ is transitive and nontrivial. A
fiberwise-linear observable (an assignment of one value per sheet) descends
to a single-valued function on $U$ if and only if it is
monodromy-invariant; the flat sections of the permutation local system
$\pi_*\C$ are the constants, and the flat endomorphisms form the
two-dimensional commutative algebra spanned by the identity and the
all-ones matrix.
\end{proposition}

\begin{proof}
Transitivity was noted above; if the monodromy were trivial the connected
total space would force $d=1$, excluded by Theorem~\ref{thm:floor}. Flat
sections are invariant vectors; a transitive permutation representation has
invariant vectors only along $\mathbf 1$, and its commutant is spanned by
$I$ and the all-ones matrix $J$.
\end{proof}

\begin{remark}\label{rem:preobservable}
Thus every counterexample possesses fiber data --- e.g.\ the nomination of
one branch --- that is locally meaningful and fully determining, yet
provably admits no globally single-valued description. Only the invariant
(symmetric) functions of the fiber descend to the base.
\end{remark}

\section{The necessary dyad}\label{sec:dyad}

From here on $F$ is a \emph{minimal} counterexample: $d=3$, so
$\Mon(\pi)=S_3$ by Theorem~\ref{thm:floor}.

The sign character $\eps\colon S_3\to\{\pm1\}$ is nontrivial on the
monodromy; its kernel defines a connected double cover
\[
U_2\longrightarrow U,
\]
the \emph{orientation cover}, on which the residual monodromy is
$A_3\cong\Z/3$. Invariantly, $U_2$ is cut out by a square root of the
discriminant class of the rank-three \'etale algebra $\pi_*\mathcal O$;
its nontriviality is exactly statement (iv) of Theorem~\ref{thm:floor}. We
call this necessary $\Z/2$ of structure \emph{the dyad}.

Let $P\to U$ be the $S_3$-Galois closure of the cover (the torsor of
ordered labelings of the three sheets). The subgroup lattice of $S_3$
provides exactly two complementary reductions:
\begin{itemize}
\item \textbf{Nomination} of a branch: the stabilizer $S_2$ of one letter
--- a \emph{non-normal} subgroup; $X_3:=P/S_2$ recovers the original
three-sheet cover.
\item \textbf{Orientation}: the kernel $A_3$ of $\eps$ --- the unique
proper \emph{normal} subgroup; $X_2:=P/A_3$ recovers the double cover
$U_2$.
\end{itemize}

\begin{proposition}[Frame factorization]\label{prop:frame}
$S_2\cap A_3=1$ and $S_2A_3=S_3$; consequently the natural map
\[
P\;\longrightarrow\;X_3\times_U X_2
\]
is an isomorphism: an ordered labeling of the three inverse branches is
exactly one nomination together with one orientation. Moreover the
determinants of the two permutation representations both equal $\eps$, so
the rank-five local system $H_2\oplus H_3$ associated to the two quotients
has trivial determinant.
\end{proposition}

\begin{proof}
Fiberwise the map is $g\mapsto(gS_2,\,gA_3)$, injective since
$S_2\cap A_3=1$ and surjective since $|S_2|\,|A_3|=|S_3|$. The permutation
action of $S_3$ on the two cosets of $A_3$ is by $\eps$, and the
determinant of the standard permutation action on three letters is $\eps$
as well; hence $\det(H_2\oplus H_3)=\eps\otimes\eps$ is trivial.
\end{proof}

\section{The coupling theorems}\label{sec:coupling}

\begin{theorem}[Uniqueness of the intermediate]\label{thm:T1}
The normal subgroups of $S_3$ are $1$, $A_3$, and $S_3$. Consequently any
resolution of the ambiguity by descent steps (connected intermediate covers
of $U$ that are Galois over $U$) has exactly one nontrivial intermediate
stage: the orientation double cover. Every Galois resolution factors
through the death of the dyad.
\end{theorem}

\begin{proof}
Connected Galois intermediates of the closure correspond to normal
subgroups; the lattice is as stated.
\end{proof}

\begin{theorem}[No orientation over the base]\label{thm:T2}
The two cyclic orderings of the three sheets --- equivalently the two
$3$-cycles $c$ and $c^{-1}=c^2$ in $S_3$ --- are conjugate in $S_3$ (by any
transposition). Hence no monodromy-invariant distinction between the two
senses of cycling exists over $U$: the unresolved ambiguity admits no
orientation.
\end{theorem}

\begin{theorem}[Orientation is born on the double cover]\label{thm:T3}
On $U_2$ the residual monodromy is the abelian group $A_3$, in which $c$
and $c^{2}$ are distinct conjugacy classes, separated by the two nontrivial
characters $\chi(c)=\omega$, $\bar\chi(c)=\omega^{2}$
\textup{(}$\omega=e^{2\pi i/3}$\textup{)}. The two cyclic orientation
classes are therefore invariantly distinguishable exactly after the dyad is
trivialized.
\end{theorem}

\begin{theorem}[The two parities]\label{thm:T3prime}
The two orientation classes of Theorem~\ref{thm:T3} are exchanged by:
\begin{enumerate}[label=\textup{(\roman*)}]
\item the deck involution of $U_2\to U$, which acts on the residual
$A_3$-data by conjugation by \textup{(}a lift of\textup{)} a transposition,
sending $c\mapsto c^{-1}$; and
\item when $F$ has real coefficients \textup{(}as the known examples
do\textup{)}, complex conjugation $\sigma$ of the base: $\sigma$ preserves
$U$ and carries each meridian of the wall to the inverse of a meridian,
transporting the monodromy representation to its inverse-composed twin and
hence exchanging $\chi$ with $\bar\chi$.
\end{enumerate}
\end{theorem}

\begin{proof}
(i) is the computation $tct^{-1}=c^{-1}$ for any transposition $t$, which
is the outer action of $S_3/A_3$ on $A_3$. For (ii): with real
coefficients the defining equations of $\SF$ are real, so $\sigma(U)=U$;
$\sigma$ is anti-holomorphic, hence orientation-reversing on each complex
transverse disk to the wall, so it carries a small positively oriented
meridian to a negatively oriented one; the induced action on local
monodromies is $g\mapsto g^{-1}$ up to conjugation, which on the abelian
residual $A_3$ exchanges the two nontrivial characters.
\end{proof}

\begin{theorem}[Real-to-complex jump]\label{thm:T4}
Every irreducible representation of $S_3$ is of real type
\textup{(}Frobenius--Schur indicator $+1$\textup{)}; in particular the
two-dimensional standard representation $V$ carried by the sheets is real.
Its restriction to $A_3$ splits as $\chi\oplus\bar\chi$, a conjugate pair
of complex-type characters \textup{(}indicator $0$\textup{)}. The
representation type of the ambiguity jumps from real to complex exactly at
the death of the dyad.
\end{theorem}

\begin{proof}
Character theory of $S_3$ and $A_3$; the indicators are classical.
\end{proof}

\begin{corollary}\label{cor:slogan}
In a minimal Jacobian counterexample, invariant orientation data
\textup{(}informally: handedness\textup{)} exists exactly and only after
the unique normal intermediate --- the sign dyad --- is resolved; and the
two available orientations are exchanged by both parity operations of
Theorem~\ref{thm:T3prime}.
\end{corollary}

\begin{remark}[What is not selected]\label{rem:grace}
Nothing in the structure selects \emph{which} orientation class is
distinguished: any labeling of $\chi$ versus $\bar\chi$ is a convention,
precisely because both parities of Theorem~\ref{thm:T3prime} exchange them.
The existence of orientation is forced; its value is not. This is
consistent with the elementary fact that no algebraic datum over $\R$
distinguishes $i$ from $-i$.
\end{remark}

\begin{remark}[A no-go, honestly recorded]\label{rem:nogo}
The orientation datum is a flat structure with finite holonomy; the
associated line systems have trivial rational Chern character, so no
ordinary characteristic-class or twisted-index invariant distinguishes the
datum from its conjugate. Whether some defect, boundary, or spectral
invariant of the wall geometry can do so is posed as
Question~\ref{q:index}.
\end{remark}

\section{Verification on the known family}\label{sec:verification}

For the explicit counterexample of \cite{Alpoge,Counterexample2026}, the
inverse cover is governed by the cubic $2ps^3-qs^2+2s-r$ over the target
$(p,q,r)$; the manuscript proves that $\SF$ is exactly the discriminant
hypersurface. Full $S_3$ monodromy is certified by two exact facts:
irreducibility, since $\C(x,y,z)=\C(p,q,r)(s)$ has degree three
\cite[Rmk.~4.4]{Counterexample2026}; and nonsquareness of the discriminant
$\Delta$, since $\operatorname{disc}_r(\Delta)=-(12p-q^2)^3\neq0$. Both
identities, and every identity below, are verified in exact rational
arithmetic by the scripts accompanying this deposit.

Across the deformation family of \cite[Thm.~5.1]{Counterexample2026}
(parameters $\lambda,a,c$ and perturbations $H$), an exact-arithmetic sweep
of $400$ specimens found full $S_3$ in every case --- as
Theorem~\ref{thm:floor} now requires; the sweep stands as verification of
the theorem rather than as evidence.

\begin{remark}[The dyad lives at infinity]\label{rem:infinity}
A Keller map is \'etale, hence unramified: no two inverse branches ever
collide at finite distance. Root collisions of the spectral polynomial over
the wall are therefore escape events --- branches leaving every compact set
--- so the transpositions constituting the sign character are monodromy
around loss-at-infinity, supported on the Jelonek set. (Compare
\cite{Anumula}, where the same phenomenon is proved for an independent
family: ``no affine ramification occurs; fiber loss is entirely a
phenomenon at infinity.'') Since Theorem~\ref{thm:floor} forces the sign
character, the discriminant class of a minimal counterexample is never a
square; equivalently, the Jelonek set always carries exchange monodromy.
\end{remark}

\section{Questions}\label{sec:questions}

\begin{question}\label{q:index}
Does any canonical defect, boundary, or spectral invariant attached to the
wall geometry distinguish the two orientation classes --- i.e., is there a
nonzero index pairing --- or does the exchange symmetry of
Theorem~\ref{thm:T3prime} force all such invariants to vanish?
\end{question}

\begin{question}\label{q:higher}
For counterexamples of degree $d\ge4$, which transitive monodromy groups
occur, and for which of them does a unique minimal normal ``orientation''
stage exist (equivalently: which have cyclic abelianization of order two)?
Is the coupling of orientation to a unique dyadic intermediate a feature of
minimality alone?
\end{question}

\begin{question}\label{q:wall}
Can the orientation cover and the local systems of
Section~\ref{sec:dyad} be extended canonically across $\SF$ (nearby cycles,
a natural singular model), and what replaces
Theorems~\ref{thm:T3}--\ref{thm:T4} on the wall itself?
\end{question}

\section*{Provenance and responsibility}

The classical inputs are exactly as cited: Keller's birational case,
Campbell's Theorem and Corollary (primary source verified against the
statements used here), the Razar and Wright generalizations, and Jelonek's
theory of the nonproperness set. The assembly of
Theorem~\ref{thm:floor}(iv), the coupling analysis of
Section~\ref{sec:coupling}, and the verification suite were developed in an
interactive research program conducted with AI systems (Anthropic's Claude
and OpenAI's ChatGPT); the author has verified every displayed argument and
computation and bears sole responsibility for all claims. Corrections,
prior references, and attributions are welcome.

\begin{thebibliography}{99}

\bibitem{Alpoge}
L.~Alp\"oge, announcement of an explicit counterexample to the Jacobian
Conjecture, July 20, 2026.

\bibitem{Counterexample2026}
\emph{A Counterexample to the Jacobian Conjecture}, verification
manuscript, July 20, 2026. \url{https://www.ulam.ai/research/jacobian.pdf}

\bibitem{Anumula}
H.~R. Anumula, \emph{Exact fibers, image, and nonproperness of the
marked-root Keller family}, Zenodo, July 21, 2026.
\href{https://doi.org/10.5281/zenodo.21506477}{doi:10.5281/zenodo.21506477}

\bibitem{Campbell}
L.~A. Campbell, \emph{A condition for a polynomial map to be invertible},
Math. Ann. \textbf{205} (1973), 243--248.
\href{https://doi.org/10.1007/BF01349234}{doi:10.1007/BF01349234}

\bibitem{Jelonek}
Z.~Jelonek, \emph{The set of points at which a polynomial map is not
proper}, Ann. Polon. Math. \textbf{58} (1993), 259--266.

\bibitem{Keller}
O.-H. Keller, \emph{Ganze Cremona-Transformationen}, Monatsh. Math. Phys.
\textbf{47} (1939), 299--306.

\bibitem{Pinchuk}
S.~Pinchuk, \emph{A counterexample to the strong real Jacobian conjecture},
Math. Z. \textbf{217} (1994), 1--4.

\bibitem{Razar}
M.~Razar, \emph{Polynomial maps with constant Jacobian}, Israel J. Math.
\textbf{32} (1979), 97--106.

\bibitem{vdE}
A.~van den Essen, \emph{Polynomial Automorphisms and the Jacobian
Conjecture}, Progress in Mathematics \textbf{190}, Birkh\"auser, Basel,
2000.

\bibitem{Wright}
D.~Wright, \emph{On the Jacobian conjecture}, Illinois J. Math.
\textbf{25} (1981), 423--440.

\end{thebibliography}

\end{document}
